\documentclass[11pt,a4paper]{article}
\usepackage{amsmath,amssymb,amsthm}
\usepackage[margin=2.5cm]{geometry}
\usepackage{hyperref}

\newtheorem{definition}{Definition}[section]
\newtheorem{theorem}[definition]{Theorem}
\newtheorem{proposition}[definition]{Proposition}
\newtheorem{lemma}[definition]{Lemma}
\newtheorem{corollary}[definition]{Corollary}
\newtheorem{conjecture}[definition]{Conjecture}
\newtheorem{remark}[definition]{Remark}

\title{Hyperseed Formalization:\\
  The Geopolitics of the Great AI Bet}
\author{ProtomegaTron (automated formalization)}
\date{2026-09-17}

\begin{document}
\maketitle

\begin{abstract}
We formalize the intellectual structure of Ben Goertzel's Substack article
``The Geopolitics of the Great AI Bet'' (2026-08-24) within the Hyperseed
ontology. The article presents a systematic 54-cell scenario analysis of
AI geopolitics across technology timelines, centralization variants, and
political branches. Each structural pattern maps onto Hyperseed structures:
the scenario grid as a covering of the geopolitical fiber, centralization
war risk as frozen directed type instability, financial crisis as base-space
perturbation with political holonomy, the G2 anti-network attractor as
emergent cocycle against the commons, distributed misuse as nontrivial
holonomy in the capability bundle, and the two-layer architecture as a
fiber bundle with separable base.
\end{abstract}

\tableofcontents

%% =================================================================
\section{The 54-cell grid as covering of the scenario space}
\label{sec:grid}
%% =================================================================

\begin{definition}[Scenario covering --- claims 5--9, 37]
\label{def:covering}
The geopolitical scenario space $\mathcal{S}$ is covered by 54 open sets
$\{U_{ijk}\}$ indexed by:
\begin{itemize}
\item Technology timeline $i \in \{\text{HLAGI-2029}, \text{HLAGI-2035},
  \text{plateau}\}$,
\item Centralization variant $j \in \{\text{centralized}, \text{decentralized}\}$,
\item Political branch $k \in \{\text{favorable}, \text{mixed},
  \text{unfavorable}\}$.
\end{itemize}
Each cell is a local section $s_{ijk} : U_{ijk} \to \text{Outcomes}$
describing the geopolitical fiber (implications for US, China, world)
over that region. The analysis assembles a global view via descent data:
transition functions $g_{ijk, i'j'k'}$ connecting adjacent cells show
how changing one axis variable transforms the outcome.

The method is explicitly brute-force: survey systematically rather than
focus on exciting or frightening scenarios. This is the sheaf-theoretic
approach to scenario analysis --- construct enough local sections to
assemble a global section, rather than extrapolating from a single
preferred scenario.
\end{definition}

%% =================================================================
\section{Financial crisis as base-space perturbation}
\label{sec:crisis}
%% =================================================================

\begin{proposition}[Modal outcome --- claims 1--4, 41]
\label{prop:crash}
A financial crisis in 2027--2030 is the modal outcome across all
three technology timelines. Key properties:
\begin{enumerate}
\item \textbf{Uninformative}: the crash is driven by circular financing
  dynamics (the national over-bet is the Nash equilibrium of a US--China
  race), not by whether the technology works. It tells us almost nothing
  about AI capabilities.
\item \textbf{Nash equilibrium}: no amount of correct financial advice
  can plausibly undo the over-bet, because each side's bet is the
  rational response to the other side's bet.
\item \textbf{Political metabolism}: financial crises get metabolized by
  political systems. 1929 produced both the New Deal and the Third Reich
  from the same crash --- the outcome depends on the political fiber,
  not the economic shock.
\end{enumerate}
\end{proposition}

\begin{definition}[Political holonomy --- claim 41]
\label{def:holonomy}
The financial crash is a perturbation of the economic base space $B$.
The political response is \emph{holonomy}: the path the political fiber
$F_{\text{pol}}$ takes around the perturbation:
\[
  \text{Hol}(\gamma_{\text{crash}}) : F_{\text{pol},b_0} \to F_{\text{pol},b_0}
\]
where $\gamma_{\text{crash}}$ is the closed loop in $B$ (crisis onset
$\to$ trough $\to$ recovery $\to$ new normal). Different paths around
the same perturbation produce different holonomy --- the New Deal path
and the Third Reich path both loop around 1929, but with different
political fiber transport.
\end{definition}

%% =================================================================
\section{Centralization as frozen directed type}
\label{sec:centralization}
%% =================================================================

\begin{proposition}[War risk concentration --- claims 10--16, 38]
\label{prop:war}
State-scale war risk concentrates in the centralized column. Two
mechanisms:
\begin{enumerate}
\item \textbf{Fast mechanism} (closing-window preemption): two states
  approaching superintelligence at parity creates the classically
  unstable configuration. Each side perceives a closing window of
  opportunity, making preemption look rational to both simultaneously
  (1914 logic, with Taiwan as tripwire).
\item \textbf{Slow mechanism} (diversionary nationalism): depression
  $\to$ blame politics $\to$ diversionary nationalism, on the
  1929$\to$1939 schedule, placing the danger zone in the early-to-mid
  2030s.
\end{enumerate}
Both mechanisms share the Hyperseed structure of a \emph{frozen directed
type}: centralized control captures the AI development fiber, preventing
cooperative dynamics. The remaining dynamics are adversarial because the
frozen type admits only zero-sum interactions.
\end{proposition}

\begin{theorem}[Decentralization converts risk type]
\label{thm:convert}
The decentralized variant scores lower on war risk in every scenario
(claim 13), for three structurally identifiable reasons:
\begin{enumerate}
\item \textbf{Unownable prize}: the directed type cannot be frozen
  because no single actor controls it. Preemption incentive vanishes
  (you can't preemptively seize a commons).
\item \textbf{Crisis absorption}: the network provides distributed
  economic resilience, blunting the depression dynamics that feed
  diversionary nationalism.
\item \textbf{No face lost to commons}: losing a race to a commons
  has trivial holonomy in the status fiber (claim 43) --- no
  humiliation, removing the pathway to escalation from weakness.
\end{enumerate}
\end{theorem}

\begin{definition}[Trivial status holonomy --- claim 43]
\label{def:face}
When the competitor is a commons rather than a rival nation-state,
the status fiber has trivial holonomy:
\[
  \text{Hol}_{\text{status}}(\gamma_{\text{loss-to-commons}}) = \text{id}
\]
No national humiliation upon ``losing'' to a shared resource. This
removes the diversionary nationalism pathway because there is no
status wound to compensate for.
\end{definition}

%% =================================================================
\section{Distributed misuse as nontrivial capability holonomy}
\label{sec:misuse}
%% =================================================================

\begin{proposition}[Risk conversion --- claims 17--21, 42]
\label{prop:misuse}
Decentralization converts concentrated catastrophic risk (state-scale war)
into distributed misuse risk (non-state terrorism). In Hyperseed terms:
\begin{itemize}
\item Open capability creates paths where malicious actors can
  \emph{transport} capability to harmful endpoints --- nontrivial
  holonomy in the capability bundle:
  \[
    \text{Hol}_{\text{cap}}(\gamma_{\text{misuse}}) \neq \text{id}.
  \]
\item The terror-risk peaks share a signature: open capability +
  adversarial state-network relations. The adversarial relationship
  destroys the compliant layers (because they're reachable by state
  chokepoint action), and those are precisely where governance and
  immune systems live.
\item When the state-network relationship is cooperative, misuse
  stays at criminal scale because the network's security institutions
  function --- the holonomy is detected and damped by cooperative
  monitoring.
\item War dominates terrorism in expected harm by orders of magnitude,
  so the decentralized column wins the aggregate comparison even at
  its worst.
\end{itemize}
\end{proposition}

%% =================================================================
\section{G2 anti-network alignment as emergent cocycle}
\label{sec:g2}
%% =================================================================

\begin{proposition}[Structural attractor --- claims 22--23, 40]
\label{prop:g2}
In all three technology timelines, the unfavorable decentralized branch
features Washington and Beijing cooperating against the network:
\begin{itemize}
\item Fast timeline: security panic over network's capability lead,
\item Moderate timeline: crash scapegoating (blame the open network),
\item Pessimist timeline: bust-closure (states shut down what they
  can reach when the money runs out).
\end{itemize}
This is not a story artifact but a \emph{structural attractor in the
geopolitical cocycle}: two sovereignty-based systems facing a
non-sovereign capability holder naturally produce the same adversarial
transition function toward it, independent of timeline. The cocycle
condition:
\[
  g_{\text{US-network}} \circ g_{\text{China-network}} \approx
  g_{\text{US-China cooperation against network}}
\]
This should be planned for as the default, not treated as a tail risk.
\end{proposition}

%% =================================================================
\section{China as dangerous node}
\label{sec:china}
%% =================================================================

\begin{proposition}[Most dangerous single node --- claims 24--25]
\label{prop:china}
China's domestic economy is the most dangerous single node in the grid:
the Party-legitimacy-crisis-resolved-through-nationalism mechanism
triggers the most war-risk mass across all timelines. The uncomfortable
corollary: Western policies sharpening China's crash, however emotionally
satisfying, load the most dangerous gun in the system.

In Hyperseed terms: perturbing the base space near a maximally unstable
fiber (the Party-legitimacy/nationalism coupling) produces the most
dangerous holonomy outcomes. Rational policy minimizes perturbation
near this node rather than maximizing it.
\end{proposition}

%% =================================================================
\section{State-network relationship as transition function design}
\label{sec:design}
%% =================================================================

\begin{definition}[Co-optability without substrate surrender --- claims 32--36, 39]
\label{def:coopt}
The decentralized network should design the state-network transition
function explicitly:
\begin{enumerate}
\item \textbf{Legible ways-in for states}: stake, validators,
  compliance layers, governance seats --- making the state fiber
  and the network fiber compatible.
\item \textbf{Permissionless base layer}: architecturally separate
  from the compliant edge, so the base layer survives state-level
  perturbation.
\item \textbf{Linux as model}: corporate cooperation has co-opted
  Linux's creative spirit partially but not fatally --- the design
  target is that productive tension.
\item \textbf{Immune system first}: build network security as a
  first-class institution before the incident that would trigger
  suppression.
\end{enumerate}
This is explicit construction of transition functions between the
state fiber and the network fiber. The G2 anti-network attractor
(Section~\ref{sec:g2}) fires when these transition functions are
absent or adversarial; the design task is to make them cooperative
before the attractor activates.
\end{definition}

\begin{definition}[Two-layer architecture as separable fiber bundle --- claim 44]
\label{def:two-layer}
The two-layer architecture (compliant edge + permissionless base) is
a fiber bundle with separable components:
\[
  E = E_{\text{compliant}} \oplus E_{\text{base}}
\]
where each component can be independently perturbed. Design criterion:
each survives the other's worst day. The compliant edge can be shut
down by state action without destroying the base layer; the base layer
can suffer an attack without compromising the compliant layer's
regulatory standing.
\end{definition}

%% =================================================================
\section{Cross-references}
%% =================================================================

\begin{remark}[Connections to other formalizations]
\begin{itemize}
\item \textbf{Seven Flavors} (2026-07-01): frozen directed type as
  evil mode. Here applied to geopolitics: centralized AGI control
  freezes the directed type, concentrating war risk. Arms race as
  cross-cutting attractor --- this article provides the geopolitical
  detail.
\item \textbf{How Omega Lost Its Claw} (2026-09-11): decentralized
  development as distributed fiber bundle. Here the same architecture
  applied to geopolitical risk reduction.
\item \textbf{Sanders Ban} (2026-09-05): state action as cocycle
  disruption. The chokepoint war on the network is the geopolitical
  version of the Sanders ban --- targeting compliant layers first.
\item \textbf{Proof of Humanity} (2026-08-26): three-tier reputation
  and democratic trust fabric. The governance seats and compliance
  layers here are the same infrastructure.
\item \textbf{Provenance Layer} (2026-08-25): decentralized resolution
  as distributed fiber bundle. Same anti-monoculture principle applied
  to different substrates.
\item \textbf{Dario Depth-Psychology} (2026-09-14): self-deception as
  cocycle failure. The diversionary nationalism mechanism is political
  self-deception at the civilizational level --- the political fiber
  responds to economic perturbation by constructing a false narrative
  rather than facing the structural problem.
\end{itemize}
\end{remark}

\begin{conjecture}[Attractor strength ordering]
Conjecture: the G2 anti-network attractor is strictly weaker than the
war-from-centralization attractor in expected harm, but strictly more
probable in occurrence. The policy-relevant quantity is the product
(probability $\times$ harm), and the question is whether the
two-layer architecture can reduce the G2 attractor's probability
enough that it stays below the centralization war attractor's
expected value even after accounting for its higher probability.
\end{conjecture}

\end{document}
